\documentclass[12pt]{article}

% Source: Tutorial 4.pdf
% Style reference: MH5200_ProblemSheet2.tex

\usepackage[margin=1in]{geometry}
\usepackage{amsmath,amssymb,mathtools}
\usepackage{enumitem}
\usepackage{bm}
\usepackage{hyperref}
\usepackage{fancyhdr}
\usepackage{draftwatermark}
\usepackage{xcolor}

%------------- TITLE AND METADATA -------------

\newcommand{\CourseCode}{MH1101 }
\newcommand{\CourseName}{Calculus II}
\newcommand{\DocType}{Tutorial 4 (Week 5) -- Problems \& Solutions}

\title{
\CourseCode \CourseName \\[0.3em]
\large \DocType
}
\author{
  Academic Year 2025/2026, Semester 2 \\[0.25em]
  \textit{Quantitative Research Society @NTU}
}
\date{\today}

%------------- HEADER & FOOTER (for page 2 onward) -------------
\fancypagestyle{main}{
  \fancyhf{}
  \fancyhead[L]{\CourseCode \CourseName}
  \fancyhead[R]{\href{https://qrsntu.org}{QRS@NTU}}
  \fancyfoot[C]{\thepage}
  \renewcommand{\headrulewidth}{0.4pt}
  \renewcommand{\footrulewidth}{0pt}
}

%------------- SIMPLE FOOTER (page 1 only) -------------
\fancypagestyle{plainfirst}{
  \fancyhf{}
  \renewcommand{\headrulewidth}{0pt}
  \renewcommand{\footrulewidth}{0pt}
  \fancyfoot[C]{\scriptsize\textcolor{gray}{\textit{For academic reference only. This document is provided for educational purposes and does not constitute an official question paper, answer key, or any formally authorised academic material. Its content is not guaranteed to be complete or accurate.}}}
}

%------------- WATERMARK -------------
\SetWatermarkText{Unofficial — QRS@NTU — qrsntu.org}
\SetWatermarkScale{1}
\SetWatermarkLightness{0.99}

\begin{document}

%------------- COVER PAGE -------------
\maketitle
\thispagestyle{plainfirst}
\pagestyle{main}

\bigskip

%====================================================
% OVERVIEW
%====================================================

\begin{center}
\rule{0.8\textwidth}{0.4pt}\\[4pt]
\textbf{Overview of This Tutorial}\\[2pt]
\rule{0.8\textwidth}{0.4pt}
\end{center}

\noindent
This tutorial develops applications of integration to volumes of revolution and provides practice with integration by parts and reduction formulas.

\begin{itemize}[leftmargin=*]
\item Cylindrical shells for volumes of revolution about vertical and horizontal axes.
\item Interpreting definite integrals as volumes of geometric solids.
\item Volume of a pyramid with equilateral triangular base via calculus and classical geometry.
\item Integration by parts for definite and indefinite integrals, including inverse trigonometric and logarithmic functions.
\item Derivation and use of a reduction formula for \(\displaystyle \int \cos^n x\,dx\).
\end{itemize}

\bigskip

\newpage

%====================================================
% QUESTION 1
%====================================================

\section*{Question 1 (Cylindrical shells)}

\subsection*{Problem}
Use the method of cylindrical shells to find the volume generated by rotating the region bounded by the given curves about the given axis.
\begin{enumerate}[label=(\alph*)]
\item \(y = x^{2}\), \(y = 6x - 2x^{2}\); about \(x = 0\).
\item \(y = x^{3}\), \(y = 8\), \(x = 0\); about \(x = 3\).
\item \(x = 2y^{2}\), \(x = y^{2} + 1\); about \(y = -2\).
\end{enumerate}

\subsection*{Solution}

\subsubsection*{Method 1: Direct shell formulas}

\begin{enumerate}[label=(\alph*)]
\item About \(x=0\) (the \(y\)-axis) with vertical shells.

The intersection points of \(y=x^{2}\) and \(y=6x-2x^{2}\) satisfy
\[
x^{2} = 6x - 2x^{2}
\;\Rightarrow\;
3x^{2}-6x=0
\;\Rightarrow\;
3x(x-2)=0,
\]
so \(x=0\) or \(x=2\). On \([0,2]\), the upper curve is \(y=6x-2x^{2}\) and the lower curve is \(y=x^{2}\). The shell at position \(x\) has
\[
\text{radius } r(x)=x,\qquad
\text{height } h(x)=(6x-2x^{2})-x^{2}=6x-3x^{2}.
\]
Thus
\begin{align*}
V
&=2\pi\int_{0}^{2} r(x)h(x)\,dx
=2\pi\int_{0}^{2} x(6x-3x^{2})\,dx\\
&=2\pi\int_{0}^{2}(6x^{2}-3x^{3})\,dx
=2\pi\left[2x^{3}-\frac{3}{4}x^{4}\right]_{0}^{2}\\
&=2\pi\left(16-12\right)=8\pi.
\end{align*}
\[
\boxed{V=8\pi.}
\]

\item About \(x=3\) with vertical shells.

The curves \(y=x^{3}\) and \(y=8\) intersect when \(x^{3}=8\), so \(x=2\). The vertical boundaries are \(x=0\) and \(x=2\). For \(x\in[0,2]\), the shell has
\[
\text{radius } r(x)=3-x,\qquad
\text{height } h(x)=8-x^{3}.
\]
Hence
\begin{align*}
V
&=2\pi\int_{0}^{2}(3-x)(8-x^{3})\,dx\\
&=2\pi\int_{0}^{2}\left(24-3x^{3}-8x+x^{4}\right)\,dx\\
&=2\pi\left[24x-\frac{3}{4}x^{4}-4x^{2}+\frac{x^{5}}{5}\right]_{0}^{2}\\
&=2\pi\left(48-12-16+\frac{32}{5}\right)
=2\pi\left(20+\frac{32}{5}\right)
=2\pi\cdot \frac{132}{5}
=\frac{264\pi}{5}.
\end{align*}
\[
\boxed{V=\dfrac{264\pi}{5}.}
\]

\item About \(y=-2\) with horizontal shells.

The intersection points of \(x=2y^{2}\) and \(x=y^{2}+1\) satisfy
\[
2y^{2}=y^{2}+1
\;\Rightarrow\;
y^{2}=1
\;\Rightarrow\;
y=\pm 1.
\]
For \(y\in[-1,1]\), the left boundary is \(x=2y^{2}\) and the right boundary is \(x=y^{2}+1\). The shell at position \(y\) has
\[
\text{radius } r(y)=y-(-2)=y+2,\qquad
\text{height } h(y)=(y^{2}+1)-2y^{2}=1-y^{2}.
\]
Thus
\begin{align*}
V
&=2\pi\int_{-1}^{1}(y+2)(1-y^{2})\,dy\\
&=2\pi\int_{-1}^{1}\big(y+2-y^{3}-2y^{2}\big)\,dy\\
&=2\pi\left[\frac{y^{2}}{2}+2y-\frac{y^{4}}{4}-\frac{2y^{3}}{3}\right]_{-1}^{1}\\
&=2\pi\left(\frac{1}{2}+2-\frac{1}{4}-\frac{2}{3}
-\left(\frac{1}{2}-2-\frac{1}{4}+\frac{2}{3}\right)\right)\\
&=2\pi\left(\frac{11}{12}+\frac{25}{12}\right)
=2\pi\cdot\frac{36}{12}
=\frac{16\pi}{3}.
\end{align*}
\[
\boxed{V=\dfrac{16\pi}{3}.}
\]
\end{enumerate}

\subsubsection*{Method 2: Alternative setups (washers / Pappus) leading to the same volumes}

\begin{enumerate}[label=(\alph*)]
\item About \(x=0\), using washers with respect to \(y\).

For \(0\le y\le 4\), the right boundary is \(x=\sqrt{y}\), and there is no left boundary (the region touches the axis), so the washer method splits at the intersection \(y=4\). For \(4\le y\le 8\), the right boundary is given implicitly by \(y=6x-2x^{2}\), i.e.
\[
x=\frac{6\pm\sqrt{36-8y}}{4},
\]
and one must select the appropriate branch. This leads to a piecewise \(R(y)\) and nontrivial algebraic expressions. The shell method in Method 1 avoids these complications and directly yields \(V=8\pi\).

\item About \(x=3\), using Pappus's centroid theorem.

The area of the region in the \(xy\)-plane is
\[
A=\int_{0}^{2}\big(8-x^{3}\big)\,dx
=\left[8x-\frac{x^{4}}{4}\right]_{0}^{2}
=16-4=12.
\]
The \(x\)-coordinate of the centroid is
\[
\bar{x}
=\frac{1}{A}\int_{0}^{2}x\big(8-x^{3}\big)\,dx
=\frac{1}{12}\left[4x^{2}-\frac{x^{5}}{5}\right]_{0}^{2}
=\frac{1}{12}\left(16-\frac{32}{5}\right)
=\frac{48}{60}
=\frac{4}{5}.
\]
When this region is rotated about \(x=3\), the centroid travels along a circle of radius \(3-\bar{x}=3-\frac{4}{5}=\frac{11}{5}\). Pappus's theorem gives
\[
V = 2\pi(3-\bar{x})A
=2\pi\cdot\frac{11}{5}\cdot 12
=\frac{264\pi}{5},
\]
matching Method 1.

\item About \(y=-2\), using washers with respect to \(x\).

For each \(x\in[0,2]\), the \(y\)-values in the region run from
\[
y_{\text{bottom}}(x)=-\sqrt{\frac{x}{2}},\qquad
y_{\text{top}}(x)=\sqrt{x-1}
\]
on the subintervals where both curves are defined, and appropriate splitting of the integral is required. The inner and outer radii from the axis \(y=-2\) are then
\[
r_{\text{in}}=y_{\text{bottom}}+2,\qquad
r_{\text{out}}=y_{\text{top}}+2.
\]
The resulting washer integral is algebraically more involved, but it evaluates to the same numerical value
\[
V = \int \pi\big(r_{\text{out}}^{2}-r_{\text{in}}^{2}\big)\,dx
= \frac{16\pi}{3},
\]
consistent with the shell computation.
\end{enumerate}

\newpage

%====================================================
% QUESTION 2
%====================================================

\section*{Question 2 (Interpreting volume integrals)}

\subsection*{Problem}
Each integral represents the volume of a solid. Describe the solid.
\begin{enumerate}[label=(\alph*)]
\item \(\displaystyle \int_{0}^{3} 2\pi x^{5}\,dx.\)
\item \(\displaystyle \int_{0}^{\pi/2} 2\pi(x+1)\,\big(2x-\sin x\big)\,dx.\)
\end{enumerate}

\subsection*{Solution}

\subsubsection*{Method 1: Reading the integrals as shell volumes}

\begin{enumerate}[label=(\alph*)]
\item The integrand has the form \(2\pi(\text{radius})(\text{height})\) with radius \(x\) and height \(x^{4}\). This is exactly the shell formula about the \(y\)-axis. The height \(x^{4}\) is the vertical distance between the curve \(y=x^{4}\) and the \(x\)-axis \(y=0\).

Therefore the solid is obtained by rotating the region bounded by
\[
y = x^{4},\quad y = 0,\quad 0\le x\le 3,
\]
about the \(y\)-axis.

\item The integrand has the form \(2\pi(\text{radius})(\text{height})\) with radius \(x+1\) and height \(2x-\sin x\). The factor \(x+1\) indicates shells about the vertical line \(x=-1\). The height \(2x-\sin x\) is the vertical distance between the top curve \(y=2x\) and the bottom curve \(y=\sin x\) for \(x\in[0,\pi/2]\).

Thus the solid is obtained by rotating the region bounded by
\[
y = 2x,\quad y = \sin x,\quad 0\le x\le \frac{\pi}{2},
\]
about the line \(x=-1\).
\end{enumerate}

\subsubsection*{Method 2: Rewriting as washer integrals}

\begin{enumerate}[label=(\alph*)]
\item The same volume can be represented via washers about the \(y\)-axis by expressing \(x\) in terms of \(y\). The curve \(y=x^{4}\) gives \(x=y^{1/4}\), and rotating the region bounded by \(y=0\), \(y=x^{4}\), \(x=0\), \(x=3\) about the \(y\)-axis yields disks of radius \(x=y^{1/4}\). The volume can be written as
\[
V = \pi\int_{0}^{81} \big(y^{1/4}\big)^{2}\,dy
= \pi\int_{0}^{81} y^{1/2}\,dy,
\]
which is equal to \(\displaystyle \int_{0}^{3} 2\pi x^{5}\,dx\) after the substitution \(y=x^{4}\).

\item For part (b), one can also regard the solid as generated by rotating the region between \(y=\sin x\) and \(y=2x\) about the line \(x=-1\). Using washers with respect to \(y\) requires solving \(y=2x\) and \(y=\sin x\) for \(x\) as functions of \(y\), which is not explicit for \(\sin x\). Hence the shell representation
\[
V=\int_{0}^{\pi/2} 2\pi(x+1)\big(2x-\sin x\big)\,dx
\]
is the natural and simplest description.
\end{enumerate}

\newpage

%====================================================
% QUESTION 3
%====================================================

\section*{Question 3 (Volume of a pyramid)}

\subsection*{Problem}
Find the volume of a pyramid with height \(h\) and base an equilateral triangle with side \(a\).

You may use the following area formula for a triangle: if \(\theta\) is the angle between two sides of lengths \(a\) and \(b\), then the area is
\[
\frac{1}{2}ab\sin\theta.
\]
You may also use Heron's formula for the area of a triangle.

\subsection*{Solution}

\subsubsection*{Method 1: Calculus with similar cross-sections}

First find the area of the base, an equilateral triangle with side \(a\). The angle between any two sides is \(\theta=\pi/3\), so
\[
A_{\text{base}}
= \frac{1}{2}a\cdot a\cdot \sin\left(\frac{\pi}{3}\right)
= \frac{1}{2}a^{2}\cdot\frac{\sqrt{3}}{2}
= \frac{\sqrt{3}}{4}a^{2}.
\]

Place the apex directly above the centroid of the base and let \(z\) measure height from the base plane (\(z=0\) at the base, \(z=h\) at the apex). A cross-section parallel to the base at height \(z\) is an equilateral triangle similar to the base, with side length scaled by a factor \((1-\frac{z}{h})\). Its area is
\[
A(z) = \left(1-\frac{z}{h}\right)^{2} A_{\text{base}}.
\]
The volume is then
\begin{align*}
V
&=\int_{0}^{h} A(z)\,dz
= A_{\text{base}}\int_{0}^{h}\left(1-\frac{z}{h}\right)^{2}\,dz\\
&= A_{\text{base}}\cdot h\int_{0}^{1}(1-u)^{2}\,du
\quad (u=z/h)\\
&= A_{\text{base}}\cdot h\left[\frac{1}{3}\right]
= \frac{1}{3}A_{\text{base}}h
= \frac{1}{3}\cdot\frac{\sqrt{3}}{4}a^{2}h.
\end{align*}
Thus
\[
\boxed{V=\frac{\sqrt{3}}{12}\,a^{2}h.}
\]

\subsubsection*{Method 2: Classical pyramid formula and base area}

A standard geometric result for any pyramid (or cone) states that its volume is
\[
V = \frac{1}{3}\times (\text{area of base}) \times (\text{height}).
\]
Using the base area computed above,
\[
A_{\text{base}}=\frac{\sqrt{3}}{4}a^{2},
\]
gives
\[
V = \frac{1}{3}\cdot \frac{\sqrt{3}}{4}a^{2}\cdot h
= \frac{\sqrt{3}}{12}a^{2}h,
\]
which agrees with Method 1:
\[
\boxed{V=\frac{\sqrt{3}}{12}\,a^{2}h.}
\]

\newpage

%====================================================
% QUESTION 4
%====================================================

\section*{Question 4 (Integration by parts)}

\subsection*{Problem}
Evaluate the integral using integration by parts.
\begin{enumerate}[label=(\alph*)]
\item \(\displaystyle \int x\cos 5x\,dx\)
\item \(\displaystyle \int \ln \sqrt{x}\,dx\)
\item \(\displaystyle \int_{0}^{1} (x^{2}+1)e^{-x}\,dx\)
\item \(\displaystyle \int_{1}^{\sqrt{3}} \tan^{-1}(1/x)\,dx\)
\item \(\displaystyle \int_{0}^{1/2} \cos^{-1} x\,dx\)
\end{enumerate}

\subsection*{Solution}

\subsubsection*{Method 1: Systematic integration by parts}

\begin{enumerate}[label=(\alph*)]
\item Let \(u=x\), \(dv=\cos 5x\,dx\). Then \(du=dx\), \(v=\frac{1}{5}\sin 5x\). Thus
\begin{align*}
\int x\cos 5x\,dx
&= uv-\int v\,du\\
&= \frac{x}{5}\sin 5x - \int \frac{1}{5}\sin 5x\,dx\\
&= \frac{x}{5}\sin 5x + \frac{1}{25}\cos 5x + C.
\end{align*}
\[
\boxed{\displaystyle \int x\cos 5x\,dx = \frac{x}{5}\sin 5x + \frac{1}{25}\cos 5x + C.}
\]

\item Note \(\ln\sqrt{x}=\frac12\ln x\). Then
\[
\int \ln\sqrt{x}\,dx
= \frac{1}{2}\int \ln x\,dx.
\]
Let \(u=\ln x\), \(dv=dx\). Then \(du=\frac{1}{x}dx\), \(v=x\). Hence
\begin{align*}
\int \ln x\,dx
&= x\ln x-\int x\cdot \frac{1}{x}\,dx
= x\ln x - x + C.
\end{align*}
Therefore
\[
\int \ln\sqrt{x}\,dx
= \frac{1}{2}\big(x\ln x - x\big) + C
= x\ln\sqrt{x} - \frac{x}{2} + C.
\]
\[
\boxed{\displaystyle \int \ln\sqrt{x}\,dx = x\ln\sqrt{x} - \frac{x}{2} + C.}
\]

\item For the definite integral
\[
\int_{0}^{1}(x^{2}+1)e^{-x}\,dx,
\]
integrate by parts twice.

First, split:
\[
\int_{0}^{1}(x^{2}+1)e^{-x}\,dx
= \int_{0}^{1}x^{2}e^{-x}\,dx + \int_{0}^{1}e^{-x}\,dx.
\]

For \(\int_{0}^{1}x^{2}e^{-x}\,dx\), let \(u=x^{2}\), \(dv=e^{-x}dx\). Then \(du=2x\,dx\), \(v=-e^{-x}\), giving
\begin{align*}
\int_{0}^{1}x^{2}e^{-x}\,dx
&= \left[-x^{2}e^{-x}\right]_{0}^{1}+\int_{0}^{1}2x e^{-x}\,dx\\
&= -e^{-1}+2\int_{0}^{1}x e^{-x}\,dx.
\end{align*}
Now integrate \(\int_{0}^{1}x e^{-x}\,dx\) by parts with \(u=x\), \(dv=e^{-x}dx\). Then \(du=dx\), \(v=-e^{-x}\), so
\begin{align*}
\int_{0}^{1}x e^{-x}\,dx
&= \left[-x e^{-x}\right]_{0}^{1}+\int_{0}^{1}e^{-x}\,dx\\
&= -e^{-1}+\big[-e^{-x}\big]_{0}^{1}\\
&= -e^{-1}+(-e^{-1}+1)=1-2e^{-1}.
\end{align*}
Therefore
\[
\int_{0}^{1}x^{2}e^{-x}\,dx
= -e^{-1}+2(1-2e^{-1})=2-5e^{-1}.
\]
Also
\[
\int_{0}^{1}e^{-x}\,dx
=\big[-e^{-x}\big]_{0}^{1}
=1-e^{-1}.
\]
Adding gives
\[
\int_{0}^{1}(x^{2}+1)e^{-x}\,dx
=(2-5e^{-1})+(1-e^{-1})
=3-6e^{-1}.
\]
\[
\boxed{\displaystyle \int_{0}^{1}(x^{2}+1)e^{-x}\,dx = 3-6e^{-1}.}
\]

\item For
\[
\int_{1}^{\sqrt{3}}\tan^{-1}(1/x)\,dx,
\]
use integration by parts with
\[
u=\tan^{-1}(1/x),\qquad dv=dx.
\]
Then
\[
du = \frac{d}{dx}\tan^{-1}\left(\frac{1}{x}\right)dx
= -\frac{1}{x^{2}+1}\,dx,\qquad v=x.
\]
Thus
\begin{align*}
\int_{1}^{\sqrt{3}}\tan^{-1}(1/x)\,dx
&=\left.x\tan^{-1}(1/x)\right|_{1}^{\sqrt{3}}
+\int_{1}^{\sqrt{3}}\frac{x}{x^{2}+1}\,dx\\
&=\left(\sqrt{3}\tan^{-1}\frac{1}{\sqrt{3}}-1\cdot\tan^{-1}1\right)
+\frac{1}{2}\left[\ln(x^{2}+1)\right]_{1}^{\sqrt{3}}\\
&=\left(\sqrt{3}\cdot\frac{\pi}{6}-\frac{\pi}{4}\right)
+\frac{1}{2}\big(\ln 4-\ln 2\big)\\
&=\frac{\sqrt{3}\pi}{6}-\frac{\pi}{4}+\frac{1}{2}\ln 2.
\end{align*}
\[
\boxed{\displaystyle \int_{1}^{\sqrt{3}}\tan^{-1}(1/x)\,dx
= \frac{\sqrt{3}\pi}{6}-\frac{\pi}{4}+\frac{1}{2}\ln 2.}
\]

\item For
\[
\int_{0}^{1/2}\cos^{-1} x\,dx,
\]
use integration by parts with
\[
u=\cos^{-1}x,\qquad dv=dx.
\]
Then \(du=-\dfrac{1}{\sqrt{1-x^{2}}}\,dx\), \(v=x\). Hence
\begin{align*}
\int_{0}^{1/2}\cos^{-1} x\,dx
&=\left.x\cos^{-1}x\right|_{0}^{1/2}
+\int_{0}^{1/2}\frac{x}{\sqrt{1-x^{2}}}\,dx\\
&=\frac{1}{2}\cos^{-1}\left(\frac{1}{2}\right)-0
+\int_{0}^{1/2}\frac{x}{\sqrt{1-x^{2}}}\,dx\\
&=\frac{\pi}{6}
+\int_{0}^{1/2}\frac{x}{\sqrt{1-x^{2}}}\,dx.
\end{align*}
For the remaining integral, let \(u=1-x^{2}\), \(du=-2x\,dx\), so
\[
\int_{0}^{1/2}\frac{x}{\sqrt{1-x^{2}}}\,dx
= -\frac12\int_{1}^{3/4}u^{-1/2}\,du
= \frac12\int_{3/4}^{1}u^{-1/2}\,du
= \left.\sqrt{u}\right|_{3/4}^{1}
= 1-\frac{\sqrt{3}}{2}.
\]
Therefore
\[
\int_{0}^{1/2}\cos^{-1} x\,dx
= \frac{\pi}{6}+1-\frac{\sqrt{3}}{2}.
\]
Equivalently,
\[
\int_{0}^{1/2}\cos^{-1} x\,dx
= \frac{1}{6}\big(\pi+6-3\sqrt{3}\big).
\]
\[
\boxed{\displaystyle \int_{0}^{1/2}\cos^{-1} x\,dx
= \frac{\pi}{6}+1-\frac{\sqrt{3}}{2}
= \frac{1}{6}\big(\pi+6-3\sqrt{3}\big).}
\]
\end{enumerate}

\subsubsection*{Method 2: Alternative viewpoints (substitutions, identities, and checking)}

\begin{enumerate}[label=(\alph*)]
\item For \(\int x\cos 5x\,dx\), tabular integration by parts (listing derivatives of \(x\) and integrals of \(\cos 5x\)) leads immediately to the same result:
\[
\int x\cos 5x\,dx = \frac{x}{5}\sin 5x + \frac{1}{25}\cos 5x + C.
\]

\item For \(\int \ln\sqrt{x}\,dx\), using the substitution \(t=\sqrt{x}\) (so \(x=t^{2}\), \(dx=2t\,dt\)) gives
\[
\int \ln\sqrt{x}\,dx
= \int \ln t\cdot 2t\,dt,
\]
which is another straightforward integration by parts in \(t\) and yields the same antiderivative.

\item For \(\int_{0}^{1}(x^{2}+1)e^{-x}\,dx\), one may expand the integrand and recognize that repeated integration by parts effectively computes the first few terms in the Taylor series of \(e^{-x}\) integrated term-wise, providing a consistency check on the value \(3-6e^{-1}\).

\item For \(\int_{1}^{\sqrt{3}}\tan^{-1}(1/x)\,dx\), the substitution \(x=1/t\) transforms the limits \([1,\sqrt{3}]\) to \([1,1/\sqrt{3}]\) and rewrites the integral in terms of \(\tan^{-1} t\), after which integration by parts proceeds similarly and confirms the same closed form.

\item For \(\int_{0}^{1/2}\cos^{-1} x\,dx\), one can use the identity
\[
\cos^{-1}x = \frac{\pi}{2}-\sin^{-1}x
\]
to relate the integral to \(\int \sin^{-1}x\,dx\), whose standard integration-by-parts formula yields
\[
\int \sin^{-1}x\,dx = x\sin^{-1}x + \sqrt{1-x^{2}}+C,
\]
leading again to the value \(\frac{1}{6}(\pi+6-3\sqrt{3})\) on \([0,1/2]\).
\end{enumerate}

\newpage

%====================================================
% QUESTION 5
%====================================================

\section*{Question 5 (Reduction formula for \(\cos^{n}x\))}

\subsection*{Problem}
\begin{enumerate}[label=(\alph*)]
\item Prove the reduction formula
\[
\int \cos^{n}x\,dx
= \frac{1}{n}\cos^{n-1}x\sin x
+ \frac{n-1}{n}\int \cos^{n-2}x\,dx.
\]
\item Use part (a) to evaluate \(\displaystyle \int \cos^{2}x\,dx\).
\end{enumerate}

\subsection*{Solution}

\subsubsection*{Method 1: Integration by parts using \(\cos^{n-1}x\cos x\)}

\begin{enumerate}[label=(\alph*)]
\item Write
\[
\int \cos^{n}x\,dx
= \int \cos^{n-1}x\cos x\,dx.
\]
Take
\[
u=\cos^{n-1}x,\qquad dv=\cos x\,dx.
\]
Then
\[
du=(n-1)\cos^{n-2}x(-\sin x)\,dx = -(n-1)\cos^{n-2}x\sin x\,dx,
\quad v=\sin x.
\]
Integration by parts gives
\begin{align*}
\int \cos^{n}x\,dx
&= uv-\int v\,du\\
&= \cos^{n-1}x\sin x - \int \sin x\big(-(n-1)\cos^{n-2}x\sin x\big)\,dx\\
&= \cos^{n-1}x\sin x + (n-1)\int \cos^{n-2}x\sin^{2}x\,dx.
\end{align*}
Using \(\sin^{2}x=1-\cos^{2}x\),
\begin{align*}
\int \cos^{n}x\,dx
&= \cos^{n-1}x\sin x + (n-1)\int \cos^{n-2}x(1-\cos^{2}x)\,dx\\
&= \cos^{n-1}x\sin x + (n-1)\int \cos^{n-2}x\,dx - (n-1)\int \cos^{n}x\,dx.
\end{align*}
Bring the last term to the left:
\[
\int \cos^{n}x\,dx + (n-1)\int \cos^{n}x\,dx
= \cos^{n-1}x\sin x + (n-1)\int \cos^{n-2}x\,dx,
\]
so
\[
n\int \cos^{n}x\,dx
= \cos^{n-1}x\sin x + (n-1)\int \cos^{n-2}x\,dx.
\]
Dividing by \(n\) yields the reduction formula
\[
\boxed{\displaystyle
\int \cos^{n}x\,dx
= \frac{1}{n}\cos^{n-1}x\sin x
+ \frac{n-1}{n}\int \cos^{n-2}x\,dx.}
\]

\item Apply the formula with \(n=2\). Then
\[
\int \cos^{2}x\,dx
= \frac{1}{2}\cos^{1}x\sin x
+ \frac{1}{2}\int \cos^{0}x\,dx
= \frac{1}{2}\cos x\sin x + \frac{1}{2}\int 1\,dx.
\]
Thus
\[
\int \cos^{2}x\,dx
= \frac{1}{2}\cos x\sin x + \frac{x}{2} + C.
\]
\[
\boxed{\displaystyle \int \cos^{2}x\,dx
= \frac{1}{2}\cos x\sin x + \frac{x}{2} + C.}
\]
\end{enumerate}

\subsubsection*{Method 2: Alternative derivations and the special case \(n=2\)}

\begin{enumerate}[label=(\alph*)]
\item For the reduction formula, an equivalent approach starts by writing
\[
\int \cos^{n}x\,dx
= \int \cos^{n-2}x(1-\sin^{2}x)\,dx
= \int \cos^{n-2}x\,dx - \int \cos^{n-2}x\sin^{2}x\,dx
\]
and then integrating the second term by parts with \(u=\sin x\) and \(dv=\cos^{n-2}x\sin x\,dx\). The resulting algebra rearranges to the same reduction formula as in Method 1.

\item For \(\int \cos^{2}x\,dx\), a widely used alternative is the double-angle identity
\[
\cos^{2}x = \frac{1+\cos 2x}{2}.
\]
Then
\begin{align*}
\int \cos^{2}x\,dx
&= \frac{1}{2}\int (1+\cos 2x)\,dx
= \frac{1}{2}\left(x+\frac{1}{2}\sin 2x\right)+C\\
&= \frac{x}{2}+\frac{1}{4}(2\sin x\cos x)+C
= \frac{1}{2}\cos x\sin x + \frac{x}{2} + C,
\end{align*}
which matches the expression obtained from the reduction formula.
\end{enumerate}

\end{document}